\subsection{A construction at a fixed limiting parameter}

Use \eqref{eq:slice}, the organizing family \eqref{intro:eq-i21-family},
and its positive speed \(f_c\) in \eqref{intro:eq-parabola-speed}.
Both infinite ends have endpoint \(q=(0,1,0)\).

\begin{theorem}\label{p4:fixed-c}
There is $K_0>0$ such that, for each fixed $c\in(1,1+K_0)$, there is
a continuous family $F_{c,\rho}$ of fields \eqref{eq:slice},
defined for all sufficiently small $\rho>0$, with
$F_{c,\rho}\to H_c$ in quadratic coefficients.
Each $F_{c,\rho}$ has three geometrically distinct hyperbolic periodic
orbits $\gamma_{j,c,\rho}$, $1\le j\le3$, and
\[
 \max_{1\le j\le3}d_{\rm H}(\iota_0(\gamma_{j,c,\rho}),\Gamma_c)
 \longrightarrow0\qquad(\rho\downarrow0).
\]
The coefficient $C_{c,\rho}$ converges to $c$.
\end{theorem}

All estimates below are uniform on the indicated compact phase and
coefficient sets. After $c$ is fixed, their constants may depend on $c$.
Every phase derivative keeps the coefficients fixed.

\subsubsection{A cubic action map with three transverse roots}

Set $k=3/2$ and
\[
 P_\eps(\xi)=(\xi+1)(\xi-\eps)(\xi-k\eps).
\]
For $0<r<1$ let $u_\eps(r)\in(0,\eps)$ be determined by
\begin{equation}\label{p4:action}
 \int_0^{u_\eps(r)}\frac{\xi}{P_\eps(\xi)}\dd\xi
 =\int_0^r\frac{s}{P_\eps(-s)}\dd s.
\end{equation}
Both integrals increase strictly and the left diverges at $u=\eps$,
giving a unique smooth $u_\eps$ with
\begin{equation}\label{p4:uprime}
 u_\eps'(r)=
 \frac{rP_\eps(u_\eps(r))}{u_\eps(r)P_\eps(-r)}>0.
\end{equation}
Define $S_\eps$ by
\[
 b_\eps(r)=-u_\eps(r)^2/2,\qquad
 S_\eps(b_\eps(r))=-r^2/2.
\]
On compact subintervals of $0<r<1$ this defines a smooth increasing
map in the phase $b$.

\begin{lemma}\label{p4:model}
There are a compact interval $J\Subset(0,1)$, three disjoint closed
subintervals $U_j\subset\operatorname{int}J$, and $\eps_*>0$ with
the following properties. Put
\[
 m_\eps=2\eps^{5/2},\qquad
 c_\eps=\frac{\eps^2}{2}-\frac{101}{100}\eps^{5/2}.
\]
For $0<\eps<\eps_*$ the function
\[
 Q_\eps(r)=\eps^{-5/2}
 \{c_\eps+b_\eps(r)-m_\eps S_\eps(b_\eps(r))\}
\]
has one simple zero in each $U_j$. Its endpoint signs and derivative
lower bounds on these intervals can be chosen independently of $\eps$.
Moreover,
\begin{equation}\label{p4:model-error}
 \left\lVert Q_\eps-\Phi\right\rVert_{C^1(J)}
 \le C_J(\sqrt\eps+\eps|\log\eps|),\qquad
 \Phi(r)=\sqrt{\frac{1-r}{8r}}+r^2-\frac{101}{100}.
\end{equation}
\end{lemma}

\phantomsection\label{app:lower-technical-proofs}
\phantomsection\label{app:lower-model-proof}

\begin{proof}[Proof of Lemma~\ref{p4:model}]
The partial fractions of $\xi/P_\eps(\xi)$ have residues
\[
 a_N=-\frac1{(1+\eps)(1+k\eps)},\quad
 a_P=-\frac2{1+\eps},\quad
 a_B=\frac3{1+k\eps},\quad a_N+a_P+a_B=0.
\]
The action difference in \eqref{p4:action} is therefore
\begin{equation}\label{p4:action-difference}
 a_N\log\frac{1+u}{1-r}
 +a_P\log\frac{\eps-u}{r+\eps}
 +a_B\log\frac{k\eps-u}{r+k\eps}.
\end{equation}
Put $\eps=h^2$, $u=h^2-dh^3$ in \eqref{p4:action-difference}.
The coefficient of $\log h$ is exactly
$3a_P+2a_B=-3h^2/((1+h^2)(1+kh^2))$; all remaining terms extend
smoothly to $h=0$, with limit $\log((1-r)/(8r))-2\log d$.
On positive compact $(d,r)$ rectangles the full $C^1_{d,r}$ error is
$O(h+h^2|\log h|)$. The limiting $d$ derivative $-2/d$ and the
implicit function theorem give
\begin{equation}\label{p4:d-error}
 d_\eps(r)=\frac{\eps-u_\eps(r)}{\eps^{3/2}},\qquad
 \left\lVert d_\eps-d_0\right\rVert_{C^1(J)}
 \le C_J(\sqrt\eps+\eps|\log\eps|),\quad
 d_0(r)=\sqrt{\frac{1-r}{8r}}.
\end{equation}
Uniqueness in \eqref{p4:action} identifies this local solution with
$u_\eps$.

The center terms cancel exactly:
\begin{equation}\label{p4:center-cancellation}
 Q_\eps(r)=d_\eps(r)-\frac{\sqrt\eps}{2}d_\eps(r)^2
               +r^2-\frac{101}{100}.
\end{equation}
This and \eqref{p4:d-error} prove \eqref{p4:model-error}.

The values at $r=1/20,1/2,49/50,999/1000$ have signs $+,-,+,-$,
as follows by evaluating
\[
 \frac{1-r}{8r}-\left(\frac{101}{100}-r^2\right)^2
 =-\frac{10000r^5-20200r^3+11451r-1250}{10000r}.
\]
Furthermore
\[
 \Phi'(r)=2r-\frac1{4\sqrt2\,r^{3/2}\sqrt{1-r}}.
\]
Its zeros are exactly the solutions of $r^5(1-r)=1/128$.
The left side increases on $(0,5/6)$ and decreases on $(5/6,1)$;
its maximum is $3125/46656>1/128$.
Thus $\Phi'$ has exactly two zeros. The alternating signs and Rolle's
theorem give exactly three roots of $\Phi$, all simple: a nonsimple
root and the two intervening Rolle points would give a third root
of $\Phi'$.
Choose $J$ containing the four test points and choose disjoint
neighborhoods $U_j$ of the three roots, indexed in increasing $r$ order,
with opposite endpoint
signs and $|\Phi'|\ge2\kappa>0$ on their union.
Reduce $\eps_*$ so that \eqref{p4:model-error} preserves these
signs and gives $|Q_\eps'|\ge\kappa$ there.
\end{proof}

Figure~\ref{fig:lower-action-profile} shows the model profile and
resolves its two roots near \(r=1\).
% Computed scalar profile used by the lower action construction.
% Native TikZ; the preamble loads tikz (no PGFPlots or external data).
% Curves sample Phi(r)=sqrt((1-r)/(8r))+r^2-101/100 directly.
% Main axes: 0.05 <= r <= 1.02, -0.45 <= Phi <= 0.15;
% the graph itself is sampled only for 0.05 <= r < 1.
% Zoom window: 0.965 <= r < 1, -0.012 <= Phi <= 0.007.
% Sampling uses a unit-interval parameter; straight segments avoid spline
% overshoot. PGF floating-point evaluation resolves the shallow right peak.
% Root markers are numerical locations (bisection residual < 5e-15):
% 0.1115973705382644, 0.9786836116269868, 0.9987531118881820.
\begin{figure}[htb]
\centering
\usepgflibrary{fpu}
\pgfmathdeclarefunction{loweractionphi}{1}{%
  \begingroup
  \pgfkeys{/pgf/fpu=true,/pgf/fpu/output format=fixed}%
  \pgfmathparse{sqrt((1-#1)/(8*#1))+#1*#1-1.01}%
  \pgfmathsmuggle\pgfmathresult
  \endgroup
}
\pgfmathdeclarefunction{loweractionphidelta}{1}{%
  \begingroup
  \pgfkeys{/pgf/fpu=true,/pgf/fpu/output format=fixed}%
  \pgfmathparse{sqrt(#1/(8*(1-#1)))-2*#1+#1*#1-.01}%
  \pgfmathsmuggle\pgfmathresult
  \endgroup
}
\begin{tikzpicture}[x=1cm,y=1cm,font=\small]
  % Panel (a): the full profile, with the divergent left tail clipped.
  \begin{scope}
    \draw[line width=.3pt] (0,0) rectangle (7.5,3.5);
    \draw[densely dashed,line width=.35pt] (0,2.625)--(7.5,2.625);
    \foreach \r/\lab in {.1/{0.1},.3/{0.3},.5/{0.5},.7/{0.7},.9/{0.9},1/{1}} {
      \pgfmathsetmacro{\xx}{(\r-.05)/.97*7.5}
      \draw[line width=.3pt] (\xx,0)--(\xx,-.065);
      \node[below,inner sep=2pt] at (\xx,-.065) {$\lab$};
    }
    \foreach \v/\lab in {-.4/{-0.4},-.2/{-0.2},0/{0},.1/{0.1}} {
      \pgfmathsetmacro{\yy}{(\v+.45)/.60*3.5}
      \draw[line width=.3pt] (0,\yy)--(-.065,\yy);
      \node[left,inner sep=2pt] at (-.065,\yy) {$\lab$};
    }
    \begin{scope}
      \clip (0,0) rectangle (7.5,3.5);
      \draw[line width=.75pt]
        plot[domain=0:1,samples=260,variable=\u]
        ({.90*\u/.97*7.5},
         {(loweractionphi(.05+.90*\u)+.45)/.60*3.5})
        plot[domain=0:1,samples=320,variable=\u]
        ({(.90+.04999*\u)/.97*7.5},
         {(loweractionphidelta(.05-.04999*\u)+.45)/.60*3.5});
      \foreach \r in {.1115973705382644,.9786836116269868,.9987531118881820} {
        \fill ({(\r-.05)/.97*7.5},2.625) circle[radius=1.25pt];
      }
    \end{scope}
    \node[above,inner sep=2pt] at (0,3.5) {$\Phi(r)$};
    \node[below,inner sep=2pt] at (7.5,-.42) {$r$};
    \node at (3.75,-.74) {(a) The action profile};
  \end{scope}
  % Panel (b): a separate linear scale resolves both right-hand zeros.
  \begin{scope}[shift={(9.3,0)}]
    \draw[line width=.3pt] (0,0) rectangle (5.2,3.5);
    \draw[densely dashed,line width=.35pt]
      (0,{.012/.019*3.5})--(5.2,{.012/.019*3.5});
    \foreach \r/\lab in {.97/{0.97},.98/{0.98},.99/{0.99},1/{1}} {
      \pgfmathsetmacro{\xx}{(\r-.965)/.035*5.2}
      \draw[line width=.3pt] (\xx,0)--(\xx,-.065);
      \node[below,inner sep=2pt] at (\xx,-.065) {$\lab$};
    }
    \foreach \v/\lab in {-.01/{-0.010},-.005/{-0.005},0/{0},.005/{0.005}} {
      \pgfmathsetmacro{\yy}{(\v+.012)/.019*3.5}
      \draw[line width=.3pt] (0,\yy)--(-.065,\yy);
      \node[left,inner sep=2pt] at (-.065,\yy) {$\lab$};
    }
    \begin{scope}
      \clip (0,0) rectangle (5.2,3.5);
      \draw[line width=.75pt]
        plot[domain=0:1,samples=360,variable=\u]
        ({.03499*\u/.035*5.2},
         {(loweractionphidelta(.035-.03499*\u)+.012)/.019*3.5});
      \foreach \r in {.9786836116269868,.9987531118881820} {
        \fill ({(\r-.965)/.035*5.2},{.012/.019*3.5})
          circle[radius=1.25pt];
      }
    \end{scope}
    \node[above,inner sep=2pt] at (0,3.5) {$\Phi(r)$};
    \node[below,inner sep=2pt] at (5.2,-.42) {$r$};
    \node at (2.6,-.74) {(b) Enlargement near $r=1$};
  \end{scope}
\end{tikzpicture}
\caption{Computed graph of
$\Phi(r)=\sqrt{(1-r)/(8r)}+r^2-101/100$ on $0<r<1$.
Panel~(a) clips the divergent left-end portion as $r\downarrow0$.
The dots mark numerically computed zeros; panel~(b) resolves the two near $r=1$.
The three simple zeros are established analytically in
Lemma~\ref{p4:model}.}
\label{fig:lower-action-profile}
\end{figure}


The lower construction balances the vanishing outer gain
$m_\eps=2\eps^{5/2}$ with the central slope
$S_\eps'(b_\eps(r))\asymp\eps^{-5/2}$ on the chosen phase interval,
as follows from \eqref{p4:d-error}. Their product remains finite and
positive; the calibration below fixes this gain and the return center
in one field.

\subsubsection{The actual passage through infinity}

Choose $K_0>0$ small enough that, for $1<c<1+K_0$,
\begin{equation}\label{p4:epsilon-c}
 \eps=\eps(c)=
 \left(\frac12e^{-4\pi/\sqrt{c-1}}\right)^{2/5}<\eps_*.
\end{equation}
Fix such a $c$ and this $\eps$ for the remainder of the construction.
Set
\begin{equation}\label{p4:path}
\begin{gathered}
 \beta=1-(1+k)\eps,\qquad \gamma=\beta\rho/3,\\
 \alpha_0=-k\eps^2\rho^3,\qquad
 \alpha_1=((1+k)\eps-k\eps^2)\rho^2,\\
 \mu_3=-\gamma-\gamma^2,\quad
 \mu_2=\alpha_1+3\gamma^2,\quad
 \mu_1=\alpha_0-\alpha_1\gamma-2\gamma^3.
\end{gathered}
\end{equation}
The coefficients $\mu$ are fixed by $c,\rho$; $A,C$ will be chosen
near $(1,c)$. In the physical field \eqref{eq:raw}, substitute
$v=\gamma+\rho\xi$, $z=\rho^2\eta$ and the positive time
$t_{\rm new}=\rho t_{\rm raw}$, giving exactly
\begin{equation}\label{p4:central}
\begin{split}
 \xi'&=-\rho P_\eps(\xi)-\ell_\rho(\xi)\eta,\\
 \eta'&=-\eta\{\xi B_\rho(\xi)+\zeta_\rho(\xi)\eta\},\\
 \ell_\rho&=1-A(\gamma+\rho\xi)+C(\gamma+\rho\xi)^2,\\
 B_\rho&=1+2\gamma+\rho\xi,\qquad
 \zeta_\rho=C\rho(\gamma+\rho\xi).
\end{split}
\end{equation}

Enlarge $J\Subset(0,1)$; its entry and exit feet remain separated from
$0$ and the three zeros of $P_\eps$. Choose compact phase intervals $I_+,I_-$
containing their images $-u^2/2,-r^2/2$, with
$b_\eps(J)\Subset\operatorname{int}I_+$ and
$S_\eps(I_+)\Subset\operatorname{int}I_-$.
Choose $L$ large enough that
\begin{equation}\label{p4:L-window}
 3/8<1/2+b/L^2<5/8
\end{equation}
on both phase intervals and on intervals connecting them to $b=0$.
On $\xi=\pm L$ use the same literal coordinate $b=\eta-L^2/2$.
The tail reference uses $b=0$.

Use the same lip height $h_\rho=(k\eps^2\rho)^{12}$ on both sides.
Between the lips use $Q=\rho\log(\eta/h_\rho)$.

\begin{lemma}\label{p4:core}
On the enlarged input interval $I_+$, the positive-time passage
of \eqref{p4:central} from $\xi=L$ to $\xi=-L$ exists for small
$\rho>0$. In the literal phase $b$ its map satisfies
\begin{equation}\label{p4:core-limit}
 \|S_{\eps,\rho;A,C}-S_\eps\|_{C^1_b}
 \le C_\eps\rho(1+|\log\rho|)\longrightarrow0,
\end{equation}
uniformly for $A,C$ in a sufficiently small fixed rectangle about
$(1,c)$, and the entire passage has $\xi'<0$.
\end{lemma}

\phantomsection\label{app:lower-core-proof}

\begin{proof}[Proof of Lemma~\ref{p4:core}]
Use the same lip height $h_\rho=(k\eps^2\rho)^{12}$ on both sides.
On each outer part put $H=\eta-\xi^2/2$ and use $\log\eta$ as a clock.
Direct differentiation gives
\begin{equation}\label{p4:fast-clock}
 \frac{\dd H}{\dd\log\eta}
 =-\frac{\rho\xi P_\eps+\xi\eta(\ell_\rho-B_\rho)
                    -\zeta_\rho\eta^2}
             {\xi B_\rho+\zeta_\rho\eta}.
\end{equation}
For the selected negative $H$ windows,
$\xi=\pm\sqrt{2(\eta-H)}$ stays a positive distance from zero.
On their fixed enlarged neighborhoods the denominator in
\eqref{p4:fast-clock} has the appropriate nonzero sign, while its
right side and the mixed derivatives needed for a first phase
variation are $O_\eps(\rho)$. Its clock length is
$O_\eps(1+|\log\rho|)$.
Integration and the scalar variational equation therefore give
\begin{equation}\label{p4:fast-error}
 |H_{\rm out}-H_{\rm in}|+
 |\partial_{H_{\rm in}}H_{\rm out}-1|
 \le C_\eps\rho(1+|\log\rho|).
\end{equation}
Applying this estimate up to a first exit from the enlarged windows
excludes such an exit for small $\rho$.

Starting at a positive lip foot $u$ and integrating outward in
reverse time gives at $\xi=L$
$b=-u^2/2+O_\eps(\rho(1+|\log\rho|))$ in $C^1_u$.
Its derivative is $-u+o(1)$, bounded away from zero.
Thus it can be inverted on the entire enlarged entry interval.
The negative lip gives in positive time
$b=-r^2/2+O_\eps(\rho(1+|\log\rho|))$ in $C^1_r$.
Near each foot, $P_\eps>0$. Away from its small foot neighborhood,
the limiting height $(\xi^2-u^2)/2$ or $(\xi^2-r^2)/2$
has a positive lower bound. This lower bound dominates the bounded
term $\rho P_\eps$, so $\rho P_\eps+\ell_\rho\eta>0$
throughout both fast parts. The hits of $\xi=\pm L$ are unique
and transverse.

Between the lips use $Q=\rho\log(\eta/h_\rho)$. In the region
$Q\le0$ the exact graph equation is
\begin{equation}\label{p4:thin}
 Q_\xi=\frac{\xi B_\rho+\zeta_\rho\eta}
                 {P_\eps+\ell_\rho\eta/\rho},
 \qquad \eta=h_\rho e^{Q/\rho}.
\end{equation}
The compact traversed $\xi$ interval, with its exit buffer, has
$P_\eps\ge p_\eps>0$.
Replace $e^{Q/\rho}$ for $Q>0$ by a bounded smooth function of
$Q/\rho$ equal to the exponential for $Q\le0$.
This defines an open scalar ODE domain and changes no solution
before its first return to the lip.
For total derivative order at most three in $Q,\xi,A,C$, the
right side differs from $\xi/P_\eps$ by at most
\begin{equation}\label{p4:thin-error}
 C_\eps\{\rho+\rho^{11-j}\},\qquad 0\le j\le3,
\end{equation}
where $j$ is the number of $Q$ derivatives. Indeed,
$h_\rho/\rho=O_\eps(\rho^{11})$, and each $Q$ derivative
costs one power of $\rho$.
The scalar integral equation and its first variational equation converge
on the fixed $\xi$ interval.

The limiting solution starting at $u$ is
$Q_0(\xi)=\int_u^\xi s/P_\eps(s)\dd s$.
It is strictly negative before its negative exit $-r$, except at
the entrance, and $Q_0'(-r)=-r/P_\eps(-r)<0$.
Uniform signs of $Q_\xi$ near entrance and exit, and a strictly
negative upper bound for $Q_0$ between them, exclude an early first
return and persist under \eqref{p4:thin-error}.
Between two buffered negative exit cuts with opposite limiting signs,
the implicit function theorem gives a unique first lip hit. The thin
field and first variations have $O_\eps(\rho)$ error and the exit has
a fixed inverse bound. Composition with \eqref{p4:fast-error} proves
\eqref{p4:core-limit}; the denominator signs and $P_\eps>0$ on the
thin part give $\xi'<0$ everywhere.
\end{proof}


\subsubsection{The two parabolic tails}

Choose a small raw half-width $d>0$. Write $q=|v-\gamma|$,
$v=\gamma+\sigma q$ with $\sigma=\pm1$, and
$\widetilde P_\rho(x)=\rho^3P_\eps(x/\rho)$.
The exact radial graph equation is
\begin{equation}\label{p4:tail}
 z_q=F_{\sigma,\rho}(q,z)=
 \frac{z\{q(1+2\gamma+\sigma q)
                  +\sigma C(\gamma+\sigma q)z\}}
 {\widetilde P_\rho(\sigma q)+
       [1-A(\gamma+\sigma q)+C(\gamma+\sigma q)^2]z}.
\end{equation}
Its initial condition at $q=\rho L$ is
$z=\rho^2(L^2/2+b)$.
The negative tail is traversed outward in positive time; the positive
tail is constructed outward in reverse time.

\begin{lemma}\label{p4:tails}
For the enlarged phase windows the tails reach the raw ports
$v=\gamma\pm d$ uniquely. Their maps $T_{\sigma,\rho}(b)$ have the
decomposition
\begin{equation}\label{p4:tail-decomposition}
 T_{\sigma,\rho}(b)=a_{\sigma,\rho}
             +c_{\sigma,\rho}\phi_{\sigma,\rho}(b),\qquad
 a_{\sigma,\rho}=T_{\sigma,\rho}(0),\quad
 c_{\sigma,\rho}=\partial_bT_{\sigma,\rho}(0)>0.
\end{equation}
Uniformly on a fixed small $A,C$ rectangle,
\begin{equation}\label{p4:tail-limits}
\begin{gathered}
 a_{\sigma,\rho}\to a_{\sigma,0},\qquad
 c_{\sigma,\rho}/\rho^2\to\ell_\sigma>0
                  \quad\text{in }C^1_{A,C},\\
 \phi_{\sigma,\rho}\to\operatorname{id}
                  \quad\text{in }C^1_b.
\end{gathered}
\end{equation}
On every fixed $q$ interval bounded away from zero, the entire
tail family converges to its limiting parabolic graph.
\end{lemma}

\phantomsection\label{app:lower-tails-proof}

\begin{proof}[Proof of Lemma~\ref{p4:tails}]
Since $\gamma=O_\eps(\rho)$ and $q\ge\rho L$,
\[
 \widetilde P_\rho(\sigma q)=O_\eps(q^3+\rho q^2+\rho^2q+\rho^3),
 \qquad |\widetilde P_\rho|/q^2\le C_\eps q.
\]
For $w=z/q^2$ in $[1/4,3/4]$ the denominator in
\eqref{p4:tail} is at least a positive constant times $q^2$
after $d$ is reduced. Its logarithmic-radius equation is
$w_{\log q}=1-2w+O_\eps(q)$, with the same bound for its
first $w$ derivative. The boundary vectors point inward.
Equation \eqref{p4:L-window} starts all phases strictly inside
this band. The band is invariant, the radial speed does not vanish,
and the raw port is reached once.

Writing the right side as
$z(\beta+\zeta z)/(\widetilde P+\ell z)$, where
$\widetilde P=\widetilde P_\rho(\sigma q)$, direct differentiation gives
\[
 F_{zz}=\frac{2\widetilde P(\zeta\widetilde P-\beta\ell)}
                 {(\widetilde P+\ell z)^3}.
 \]
Together with the band bounds this yields
\begin{equation}\label{p4:tail-majorants}
\begin{split}
 |F_z|&\le C_\eps(1+\rho/q+\rho^2/q^2+\rho^3/q^3),\\
 |F_{zz}|&\le C_\eps q^{-2},\\
 |F_A|+|F_C|&\le C_\eps q^2,\qquad
 |F_{zA}|+|F_{zC}|\le C_\eps .
\end{split}
\end{equation}
The first majorant is uniformly bounded since $\rho/q\le1/L$.
For $J=\partial_bz$, $J(\rho L)=\rho^2$ and $J_q=F_zJ$.
Consequently
\begin{equation}\label{p4:J}
 c_\eps'\rho^2\le J\le C_\eps'\rho^2,\qquad
 |\partial_b\log J|
 \le\int_{\rho L}^d|F_{zz}|J\dd q\le C_\eps'\rho.
\end{equation}
This proves the last assertion of \eqref{p4:tail-limits}, since
$\phi(0)=0$ and $\phi'=J(q,b)/J(q,0)$.

At $\rho=0$ equation
\eqref{p4:tail} has a unique graph in the band with
$z_0(q)=q^2/2+O(q^3)$.
It is obtained by starting $z(a)=a^2/2$ and letting $a\downarrow0$.
The inward-band comparison gives $|z/q^2-1/2|\le Cq$
for these initial data. Any two bounded-band limiting graphs have
distance at most $Ca^2$ at $q=a$; the bound for $F_z$
propagates this distance to $q=d$, proving uniqueness.
For the actual reference $b=0$, coefficient convergence on compact
$q>0$ intervals, $z\le Cq^2$, and propagation from a small fixed $a$
give convergence to this graph.

The parameter-independent initial height gives
$z_{A,q}=(\partial_zF)z_A+F_A$ and the analogous equation for $C$.
The majorants \eqref{p4:tail-majorants} give $|z_A|+|z_C|\le Cq^3$;
their integral equations converge to the unique zero-initial-value
equations at $q=0$, proving $C^1_{A,C}$ convergence of $a_{\sigma,\rho}$.
Finally,
\[
 \log(c_{\sigma,\rho}/\rho^2)=
 \int_{\rho L}^dF_z(q,z(q,0))\dd q.
\]
Its $A,C$ derivatives have integrands
$F_{zA}+F_{zz}z_A$ and $F_{zC}+F_{zz}z_C$, bounded by
$C(1+q)$.
Pointwise convergence for $q>0$ and these integrable majorants give
the first-slope limits in \eqref{p4:tail-limits}.
The same bounds show $\ell_\sigma$ and their reciprocals are bounded
on the coefficient rectangle.
\end{proof}


\subsubsection{The finite connection and two coefficient directions}

At $\rho=0,A=1,C=c$ the two limiting tail graphs are the ends of
the organizing parabola, whose raw height is
\begin{equation}\label{p4:raw-organizing}
 z_c(v)=\frac{v^2}{1+\sqrt{1+cv^2}}.
\end{equation}
The finite arc between $v=-d$ and $v=d$ is a compact affine arc.
It includes the two crossings $Y=0,\ X=\pm\sqrt c$.
At these points
\[
 \dot X=c\pm\sqrt c>0,\qquad
 \dot Y=c\pm c^{3/2}\ne0.
\]
Along the entire arc $f_c\ge(c-1)/2$.
Choose a fixed narrow affine tube and buffered raw section windows.
After the coefficient rectangle and $\rho$ are reduced, the
finite-time flow stays in this tube and hits the positive port
transversely. It defines an increasing map $E_\rho$ from negative
to positive raw height. Its mixed derivatives through order two
converge uniformly to those of $E_0$.
The same transverse finite-time estimates give a uniform $C^3$ norm
and slope bounded away from zero.
Lemma~\ref{p4:tails} admits the port heights and all phase images to
these buffered domains before composition.

Define
\begin{equation}\label{p4:outer-data}
 \Delta_\rho=a_{+,\rho}-E_\rho(a_{-,\rho}),\qquad
 \Lambda_\rho=E_\rho'(a_{-,\rho})c_{-,\rho}/c_{+,\rho}.
\end{equation}
These functions converge in $C^1_{A,C}$ by the finite-time
estimates and the normalized slope limits in \eqref{p4:tail-limits}.
The limiting return center varies nontrivially with $A$, while $C$
changes the complete gain along the connected family $A=1$.
The triangular Jacobian below allows both conditions to be imposed
on the same field.

\begin{lemma}\label{p4:transverse-coefficients}
At $(A,C)=(1,c)$,
\begin{equation}\label{p4:Jacobian}
 \Delta_0=0,\quad \partial_C\Delta_0=0,\quad
 \partial_A\Delta_0>0,\qquad
 \Lambda_0=e^{-4\pi/\sqrt{c-1}},\quad
 \partial_C\Lambda_0=2\pi(c-1)^{-3/2}\Lambda_0>0.
\end{equation}
\end{lemma}

\phantomsection\label{app:lower-coefficients-proof}

\begin{proof}[Proof of Lemma~\ref{p4:transverse-coefficients}]
For $A=1$ the whole $C>1$ family has the invariant parabola, so
$\Delta_0(1,C)=0$.
Write $K=c-1$ and
\[
 \mathcal E_c(X)=((X+1)^2+K)^2
 \exp\left[-\frac4{\sqrt K}
                \arctan\frac{X+1}{\sqrt K}\right].
\]
For $W=Y-(X^2-c)/2$, the exact affine equation at $\mu=0,C=c$ is
\[
 \dot W=2XW+(1-A)X^2,\qquad
 \mathcal E_c'/\mathcal E_c=2X/f_c.
\]
The graph variation in the $A$ direction therefore satisfies
\[
 W_{A,X}=(2X/f_c)W_A-X^2/f_c.
\]
The tail bound $z_A=O(q^3)$ gives $W_A=O(X)$ at infinity.
Since $\mathcal E_c\asymp X^4$, the free homogeneous coefficient
vanishes, and the selected variations are
\[
 W_A^-(X)=-\mathcal E_c(X)\int_{-\infty}^X
       \frac{t^2}{f_c(t)\mathcal E_c(t)}\dd t,\qquad
 W_A^+(X)=\mathcal E_c(X)\int_X^\infty
       \frac{t^2}{f_c(t)\mathcal E_c(t)}\dd t.
\]
The positive integrand is $O(|t|^{-4})$ at both ends, so
$W_A^+-W_A^-$ is positive at the positive port. At fixed $v=d$,
\[
 W_z=\frac{d^2-z_c(d)}{z_c(d)^3}>0,
\]
conversion to raw height gives $\partial_A\Delta_0>0$.

Truncate the limiting tails at $v=\pm a$. The outer derivative is
$E_0'$ times the negative-tail first variation divided by the positive
one, tending to $\Lambda_0=E_0'\ell_-/\ell_+$ as $a\downarrow0$.
The parabola endpoints are $X=\pm X_a$, $X_a=a/z_c(a)\to\infty$,
with equal fixed-$v$ conversion factors $W_z$. Thus the derivative
is $\mathcal E_c(X_a)/\mathcal E_c(-X_a)$: its polynomial ratio tends
to one and its arctangent difference to $\pi$, giving
$\Lambda_0(1,c)=e^{-4\pi/\sqrt K}$. For every $C$ in the rectangle,
$\Lambda_0(1,C)=e^{-4\pi/\sqrt{C-1}}$, so differentiation gives
$\partial_C\Lambda_0(1,c)=2\pi K^{-3/2}e^{-4\pi/\sqrt K}$.
\end{proof}


\subsubsection{A single calibration and three simple actual roots}

Equation \eqref{p4:epsilon-c} gives
$\Lambda_0(1,c)=m_\eps$ exactly.
Consider the two equations
\begin{equation}\label{p4:calibration}
 G_\rho(A,C)=
 \bigl(\Delta_\rho-c_{+,\rho}c_\eps,\,
                         \Lambda_\rho-m_\eps\bigr)=0.
\end{equation}
By \eqref{p4:tail-limits},
$c_{+,\rho}\to0$ together with its parameter derivatives.
Thus $G_\rho\to G_0$ in $C^1_{A,C}$ and $G_0(1,c)=0$.
Lemma~\ref{p4:transverse-coefficients} gives
\[
 \det DG_0(1,c)
 =(\partial_A\Delta_0)(\partial_C\Lambda_0)>0.
\]
On a small fixed ball in the previous buffered domains with $C>1$,
$C^1$ convergence makes
$(A,C)\mapsto(A,C)-DG_0(1,c)^{-1}G_\rho(A,C)$ a contraction into
itself for small $\rho$. The resulting coefficient pair satisfies
\begin{equation}\label{p4:AC}
 (A_\rho,C_\rho)\longrightarrow(1,c).
\end{equation}
The calibrated pair is continuous in $\rho>0$, since all actual hits
are transverse.
No estimate uniform in $c\downarrow1$ is needed for this inverse.

Evaluate all maps at this one coefficient pair. Normalize
$E_\rho\circ T_{-,\rho}$ by its exact center and slope:
\[
 E_\rho(T_{-,\rho}(b))=
 E_\rho(a_{-,\rho})+
 E_\rho'(a_{-,\rho})c_{-,\rho}\widehat\phi_{-,\rho}(b).
\]
Here $\widehat\phi_{-,\rho}=N_{c_{-,\rho}}\circ\phi_{-,\rho}$, with
$\psi=E_\rho$ and $z_0=a_{-,\rho}$ in
\eqref{phys:eq-observer-normalization}. The buffered $C^3$ bounds,
$c_{-,\rho}=O(\rho^2)$, and \eqref{p4:tail-limits} give
$\widehat\phi_{-,\rho}(0)=0$ and $C^1_b$ convergence to the identity.

The return equation is
$T_{+,\rho}(b)=E_\rho(T_{-,\rho}(S_{\eps,\rho}(b)))$.
Dividing its difference by the positive $c_{+,\rho}$ and using
\eqref{p4:calibration} gives exactly
\begin{equation}\label{p4:actual-displacement}
 D_\rho(b)=c_\eps+\phi_{+,\rho}(b)
         -m_\eps\widehat\phi_{-,\rho}(S_{\eps,\rho}(b)).
\end{equation}
For $D_\eps=c_\eps+b-m_\eps S_\eps(b)$, the core bound,
\eqref{p4:J}, and the buffered finite connector give, on the common
enlarged interval and uniformly in the coefficient ball,
\begin{equation}\label{p4:normalized-error}
 \left\lVert D_\rho-D_\eps\right\rVert_{C^1_b}
       \le C_\eps\rho(1+|\log\rho|).
\end{equation}
This remains true at \eqref{p4:AC}. Keep $\eps=\eps(c)$ fixed and put
$\mathcal Q_\rho(r)=\eps^{-5/2}D_\rho(b_\eps(r))$.
The chain rule gives the explicit estimate
\begin{equation}\label{p4:actual-error}
 \left\lVert \mathcal Q_\rho-Q_\eps\right\rVert_{C^1(J)}
 \le \eps^{-5/2}\max(1,\left\lVert b_\eps'\right\rVert_{C^0(J)})
                         \left\lVert D_\rho-D_\eps\right\rVert_{C^1_b}
 \longrightarrow0.
\end{equation}
All coefficients, including $A_\rho,C_\rho$, are fixed in this
$r$ derivative. For the fixed $\eps$, choose the same $\rho$
small enough to preserve the three endpoint sign pairs and the
derivative lower bound of Lemma~\ref{p4:model}.
There is exactly one root $r_{j,\rho}$ in each $U_j$, and
$|\mathcal Q_\rho'(r_{j,\rho})|\ge\kappa/2$.

Each root closes one actual circuit. The central passage and both
tails have strictly decreasing $v$ in positive time.
For each $\rho>0$ all flight times are finite: on the thin part
$-\xi'\ge\rho p_\eps$, the fast parts are compact with nonzero
speed, and the outward tail speed is at least a positive constant
times $q^2$ with $q\ge\rho L>0$.
The finite connection lies in a compact affine tube whose spherical
distance from $q$ is positive. The inner section $\xi=L$ tends to
$q$, so for small $\rho$ it has no further hit during the circuit.
Thus the three roots give three different prime periodic orbits.

They are hyperbolic in the actual raw input coordinate.
Indeed $h(r)=T_{+,\rho}(b_\eps(r))$ is positive and
\[
 h_r=c_{+,\rho}\phi_{+,\rho}'b_\eps'(r)<0,\qquad h_r\ne0,
\]
by \eqref{p4:uprime}. If $R_\rho$ is the first return and
$F_\rho(h)=h-R_\rho(h)$, then the exact identity
\[
 F_\rho(h(r))=c_{+,\rho}\eps^{5/2}\mathcal Q_\rho(r)
\]
gives, at each root,
\begin{equation}\label{p4:physical-simple}
 F_{\rho,h}=
 \frac{\eps^{5/2}\mathcal Q_\rho'}
      {\phi_{+,\rho}'b_\eps'}\ne0.
\end{equation}
The common gain $c_{+,\rho}$ cancels; positivity of the first-return
derivative and $R_\rho'\ne1$ prove hyperbolicity.

\subsubsection{Whole-orbit convergence}

The three roots remain in the same compact phase intervals.
On the central passage $|\xi|\le L$ and $0<\eta\le C_{\eps,L}$.
At $\xi=0$ the thin orbit has $Q\le0$, hence all three satisfy
\begin{equation}\label{p4:third-height}
 0<z(\gamma)=\rho^2\eta(0)
 \le\rho^2h_\rho=(k\eps^2)^{12}\rho^{14}\longrightarrow0.
\end{equation}

Fix $0<d_0<d$. On $d_0\le|v-\gamma|\le d$, both tails converge
uniformly to their organizing ends, by Lemma~\ref{p4:tails} and
\eqref{p4:AC}. The finite connection converges uniformly, in affine
coordinates and finite time, to the corresponding organizing arc.
Each of these arcs is parameterized over a full common interval,
with endpoints converging to the corresponding organizing endpoints.
Consequently both directed Hausdorff distances tend to zero on
these pieces.

The remaining tail pieces have $z\le C|v-\gamma|^2$ and
$|v|\le d_0+O(\rho)$.
The central piece has $|v-\gamma|\le\rho L$ and $z\le C\rho^2$.
For these pieces $Y>0$, and the true spherical embedding is
\[
 \iota_0(X,Y)=\frac{(v,1,z)}{\sqrt{1+v^2+z^2}}.
\]
Their distance to $q$ is at most $C(d_0+\rho)$.
The omitted organizing ends have distance at most $Cd_0$ to $q$.
Conversely, \eqref{p4:third-height} provides an actual point tending
to $q$ on each of the three cycles, so every omitted organizing
point is within $C d_0+o(1)$ of that cycle.
First let $\rho\downarrow0$, then $d_0\downarrow0$.
This proves simultaneous whole-orbit Hausdorff convergence and
Theorem~\ref{p4:fixed-c}. The finite part containing $Y=0$
uses the affine embedding throughout.

\begin{corollary}[Stability order]\label{p4:stability-order}
For each fixed $c\in(1,1+K_0)$ and all sufficiently small $\rho>0$,
order the three constructed cycles of $F_{c,\rho}$ by their roots
$r_{1,\rho}<r_{2,\rho}<r_{3,\rho}$. The first and third cycles
are attracting, and the second is repelling.
\end{corollary}
\begin{proof}
The roots of \(\Phi\), in increasing \(r\) order, cross downward,
upward and downward. The two \(C^1\) estimates
\eqref{p4:model-error} and \eqref{p4:actual-error} preserve these
signs for \(\mathcal Q_\rho'\). Since \(b_\eps'<0\) and
\(\phi_{+,\rho}'>0\), identity \eqref{p4:physical-simple} gives
signs \(+,-,+\) for \(1-R_\rho'\). Positivity of the actual
first-return derivative proves the stated stability order.
\end{proof}

\begin{corollary}[Limits of the three return multipliers]
\label{p4:multiplier-limits}
Let $r_1^0<r_2^0<r_3^0$ be the three zeros of $\Phi$ and let
$M_j(c,\rho)>0$ be the derivative of the positive first return at
the cycle $\gamma_{j,c,\rho}$, ordered by increasing $r$.
For $\eps=\eps(c)$, let $r_{j,\eps}$ be the zero of $Q_\eps$ in
$U_j$ and put $\mu_j(c)=m_\eps S_\eps'(b_\eps(r_{j,\eps}))>0$.
Then
\begin{equation}\label{p4:multiplier-successive-limits}
 \begin{gathered}
 \lim_{\rho\downarrow0}M_j(c,\rho)=\mu_j(c),\qquad
 \lim_{c\downarrow1}\mu_j(c)=L(r_j^0),\\
 L(r)=8\sqrt2\,r^{5/2}\sqrt{1-r}.
 \end{gathered}
\end{equation}
The limiting multipliers satisfy
\begin{equation}\label{p4:multiplier-intervals}
 0<L(r_1^0),L(r_3^0)<\tfrac12,
 \qquad \tfrac32<L(r_2^0)<2.
\end{equation}
\end{corollary}
\begin{proof}
For fixed $c$ and $\eps=\eps(c)$, \eqref{p4:actual-error} gives
$r_{j,\rho}\to r_{j,\eps}$. Evaluate all maps at the calibrated
coefficients \eqref{p4:AC}. In the actual raw input coordinate
$h=T_{+,\rho}(b)$, the first return is
\[
 R_\rho=E_\rho\circ T_{-,\rho}\circ
 S_{\eps,\rho}\circ T_{+,\rho}^{-1}.
\]
At $b=b_\eps(r_{j,\rho})$, its derivative is exactly
\begin{equation}\label{p4:multiplier-exact}
 \begin{split}
 M_j(c,\rho)
 &=\frac{E_\rho'(T_{-,\rho}(S_{\eps,\rho}(b)))
          T_{-,\rho}'(S_{\eps,\rho}(b))
          S_{\eps,\rho}'(b)}{T_{+,\rho}'(b)}\\
 &=m_\eps\,
   \frac{\widehat\phi_{-,\rho}'(S_{\eps,\rho}(b))
          S_{\eps,\rho}'(b)}{\phi_{+,\rho}'(b)}>0.
 \end{split}
\end{equation}
The second equality uses the exact calibration
$\Lambda_\rho=m_\eps$. The core and tail $C^1$ limits, including
$\widehat\phi_{-,\rho}\to\operatorname{id}$, give the first limit
in \eqref{p4:multiplier-successive-limits}.

All derivatives keep the coefficients fixed. From
$u_\eps=\eps-\eps^{3/2}d_\eps$ and
$S_\eps(b_\eps(r))=-r^2/2$ we obtain
\begin{equation}\label{p4:balanced-slope-limit}
 \begin{gathered}
 b_\eps'(r)=\eps^{5/2}(1-\sqrt\eps\,d_\eps(r))d_\eps'(r),\\
 m_\eps S_\eps'(b_\eps(r))
 =\frac{-2r}{(1-\sqrt\eps\,d_\eps(r))d_\eps'(r)}
 \longrightarrow L(r)
 \quad\text{uniformly on }J.
 \end{gathered}
\end{equation}
Indeed \eqref{p4:d-error} gives $d_\eps\to d_0$ in $C^1(J)$ and
$d_0'(r)=-1/(4\sqrt2\,r^{3/2}\sqrt{1-r})$ is bounded away from
zero there. By \eqref{p4:model-error}, $r_{j,\eps}\to r_j^0$.
Together with $\eps(c)\to0$, this proves
\eqref{p4:multiplier-successive-limits}.

The bounds needed for \eqref{p4:multiplier-intervals} follow from
exact rational brackets. The sign of $\Phi(r)$ on $(0,1)$ is the
sign of
\[
 \frac{1-r}{8r}-\left(\frac{101}{100}-r^2\right)^2
 =-\frac{10000r^5-20200r^3+11451r-1250}{10000r}.
\]
Direct substitution gives the endpoint signs in the following table.
Since $L(r)^2=128r^5(1-r)$ increases up to $r=5/6$ and decreases
afterwards, rational endpoint evaluation also gives its indicated
bounds throughout each bracket. The endpoints have common denominator
$10^5$:
\[
\begin{array}{c|c|c|c}
 j&(10^5a_j,10^5b_j)&(\operatorname{sgn}\Phi(a_j),
                  \operatorname{sgn}\Phi(b_j))
       &L(r)^2\ (a_j\le r\le b_j)\\ \hline
 1&(11159,11160)&(+,-)&(1/1000,1/400)\\
 2&(97868,97869)&(-,+)&(12/5,5/2)\\
 3&(99875,99876)&(+,-)&(3/20,1/5)
\end{array}
\]
Each bracket contains $r_j^0$, since $\Phi$ has exactly three
roots. Taking positive square roots gives
\eqref{p4:multiplier-intervals}.
\end{proof}
